\blattunter{Blatt 0 · Das kannst du schon}

\begin{aufgabe}{Quadratzahlen -- im Kopf. Berechne.}
\begin{teilezwei}
\tz $6^2 = \leerfeld$ & \tz $9^2 = \leerfeld$ \\
\tz $0{,}4^2 = \leerfeld$ & \tz $(-7)^2 = \leerfeld$ \\
\tz $-7^2 = \leerfeld$ & \tz $3 \cdot 4^2 = \leerfeld$ \\
\end{teilezwei}
\end{aufgabe}

\begin{aufgabe}{Quadratzahlen -- Fehler finden. Nele rechnet:}
\rechnung{(-9)^2 &= -81}
Finde den Fehler, benenne ihn und korrigiere die Rechnung.

\begin{teile}
\teil Berechne $(-12)^2$ und $-12^2$.
\end{teile}
\end{aufgabe}

\begin{aufgabe}{Quadratwurzeln -- ziehen. Berechne im Kopf.}
\begin{teilezwei}
\tz $\sqrt{36} = \leerfeld$ & \tz $\sqrt{81} = \leerfeld$ \\
\tz $\sqrt{169} = \leerfeld$ & \\
\end{teilezwei}
Rechne mit dem Taschenrechner und runde auf zwei Stellen nach dem Komma.
\begin{teilezwei}
\tz $\sqrt{7} \approx \leerfeld$ & \tz $\sqrt{30} \approx \leerfeld$ \\
\end{teilezwei}
Gibt es die Wurzel? Wenn ja, berechne sie; wenn nein, schreib „gibt es nicht``.
\begin{teilezwei}
\tz $\sqrt{-16} = \leerfeld$ & \tz $\sqrt{0} = \leerfeld$ \\
\end{teilezwei}
\end{aufgabe}

\begin{aufgabe}{Gleichungen -- mit dem Strich lösen. Schreib die Umformung hinter den
Strich und rechne darunter weiter.}
\beispiel{x + 6 &= 10 &&\mid -6 \\ x + 6 - 6 &= 10 - 6 \\ x &= 4}
\begin{gleichungsraster}[2]
\gl{x + 7 = 15} & \gl{x - 5 = 11} \\
\gl{4x = 28} & \gl{-x = 12} \\
\gl{6x = -42} & \gl{x + 14 = 5} \\
\end{gleichungsraster}
\end{aufgabe}

\begin{aufgabe}{Lösungen -- durch Einsetzen prüfen. Setze die Zahl für $x$ ein, rechne
beide Seiten aus und kreuze an.}
\beispiel{2x - 1 &= 9 && x = 5 \\ 2 \cdot 5 - 1 &= 9 \\ 9 &= 9 && \text{(wA)}}
\begin{teile}
\teil $3x - 5 = 7$ \quad für $x = 4$ \hfill \kreuz{wahr}\kreuz{falsch}
\teil $x + 9 = 14$ \quad für $x = 6$ \hfill \kreuz{wahr}\kreuz{falsch}
\teil $x + 11 = 9$ \quad für $x = -3$ \hfill \kreuz{wahr}\kreuz{falsch}
\teil $x^2 = 25$ \quad für $x = -5$ \hfill \kreuz{wahr}\kreuz{falsch}
\teil $x^2 + 3x = 10$ \quad für $x = -2$ \hfill \kreuz{wahr}\kreuz{falsch}
\teil $x \cdot (x + 9) = -20$ \quad für $x = -4$ \hfill \kreuz{wahr}\kreuz{falsch}
\end{teile}
\end{aufgabe}

\begin{aufgabe}{Scheitelpunktform -- lesen. Lies den Scheitel ab und kreuze an, wohin
die Parabel geöffnet ist.}
\begin{teile}
\teil $p(x) = (x - 3)^2 + 2$ \quad \punktfeld \quad \kreuz{oben}\kreuz{unten}
\teil $p(x) = (x - 6)^2 + 1$ \quad \punktfeld \quad \kreuz{oben}\kreuz{unten}
\teil $p(x) = (x - 2)^2 - 7$ \quad \punktfeld \quad \kreuz{oben}\kreuz{unten}
\teil $p(x) = (x + 5)^2 + 3$ \quad \punktfeld \quad \kreuz{oben}\kreuz{unten}
\teil $p(x) = -(x - 4)^2 + 6$ \quad \punktfeld \quad \kreuz{oben}\kreuz{unten}
\teil $p(x) = x^2 - 12$ \quad \punktfeld \quad \kreuz{oben}\kreuz{unten}
\end{teile}
\end{aufgabe}

\begin{aufgabe}{Klammern -- zuerst rechnen. Rechne zuerst in der Klammer.}
\begin{teilezwei}
\tz $5 \cdot (3 + 4) = \leerfeld$ & \tz $(-6) \cdot (2 + 5) = \leerfeld$ \\
\tz $(-4) \cdot (-4 + 9) = \leerfeld$ & \tz $(-3) \cdot (-3 - 5) = \leerfeld$ \\
\tz $7 \cdot (7 - 12) = \leerfeld$ & \tz $(-8) \cdot (-8 + 8) = \leerfeld$ \\
\end{teilezwei}
\end{aufgabe}

\begin{aufgabe}{Klammern -- ausmultiplizieren und ausklammern. Multipliziere aus und
fasse zusammen.}
\begin{teile}
\teil $x \cdot (x + 6) = \leerfeld$
\teil $x \cdot (x - 9) = \leerfeld$
\teil $(x + 5)^2 = \leerfeld$
\teil $(x - 7)^2 = \leerfeld$
\teil $(x + 3) \cdot (x - 8) = \leerfeld$
\end{teile}
Klammere $x$ aus.
\begin{teile}
\teil $x^2 + 11x = \leerfeld$
\teil $x^2 - 6x = \leerfeld$
\end{teile}
\end{aufgabe}

\begin{aufgabe}{Terme -- ordnen und zusammenfassen. Ordne nach $x^2$, dann $x$, dann
Zahl und fasse zusammen.}
\begin{teile}
\teil $x^2 + 4x + 9x = \leerfeld$
\teil $3x + x^2 - 8 = \leerfeld$
\teil $x^2 - 7 + 5x + 12 = \leerfeld$
\teil $2x - x^2 + 9x = \leerfeld$
\teil $x^2 + 6x - 4 - 10x = \leerfeld$
\teil $x^2 - 9x + 2x - 6 = \leerfeld$
\end{teile}
\end{aufgabe}

\begin{aufgabe}{Punkte -- aufschreiben und ablesen. Schreib den Punkt in der Form
$P(x \mid y)$ auf.}
\begin{teile}
\teil $3$ nach rechts, $5$ nach oben \hfill \pfeld{P}
\teil $2$ nach rechts, $4$ nach unten \hfill \pfeld{P}
\teil $6$ nach links, auf der $x$-Achse \hfill \pfeld{P}
\end{teile}
Lies die drei Punkte ab.

\begin{ksys}[xmin=-6,xmax=6,ymin=-5,ymax=5,ablesen]
\punkt{-4}{3}{A}
\punkt{3}{-2}{B}
\punkt{-2}{-4}{C}
\end{ksys}

\begin{teile}
\teil \pfeld{A}
\teil \pfeld{B}
\teil \pfeld{C}
\end{teile}
\end{aufgabe}

\begin{aufgabe}{Geraden -- Wert berechnen und gleichsetzen. Berechne den Funktionswert.}
\begin{teile}
\teil $f(x) = 2x + 1$ \quad $f(3) = \leerfeld$
\teil $f(x) = -x + 8$ \quad $f(5) = \leerfeld$
\teil $f(x) = 3x - 4$ \quad $f(-2) = \leerfeld$
\end{teile}
Setze die beiden Terme gleich und berechne $x$.
\begin{gleichungsraster}[3]
\gl{2x + 1 = x + 6} & \gl{3x - 5 = -x + 11} \\
\gl{-2x + 1 = 4x - 17} & \\
\end{gleichungsraster}
Berechne zu f) auch den $y$-Wert und schreib den Schnittpunkt auf.
\begin{teile}
\teil \pfeld{S}
\end{teile}
\end{aufgabe}
